\documentclass[11pt,oneside,a4paper]{book}
\usepackage[utf8]{inputenc}
\usepackage[spanish, es-tabla, es-lcroman]{babel}  
\usepackage{calc}
\usepackage{enumerate}    
\usepackage{etoolbox}
\usepackage{hyperref}                       
\usepackage{amsfonts}
\usepackage{graphicx}
\usepackage{wrapfig}
\usepackage{float}
\usepackage{amssymb}
\usepackage{amsmath}
\usepackage{gensymb}
\usepackage{changepage}
\usepackage[nottoc]{tocbibind}
\usepackage{tikz} 
\tikzset{font=\fontfamily{phv}\selectfont}
\usepackage{pgfplots}
\usepackage{ragged2e}
\usepackage{caption}       
\usepackage[justification=justified,singlelinecheck=false]{caption}
\DeclareCaptionLabelSeparator{mysep}{.\hspace{2mm}}
\captionsetup{format=hang, labelsep=mysep}
\usepackage{tocloft}
\usepackage{cleveref}
\usepackage{tocvsec2}
\usepackage{appendix}
\usepackage{rotating}       
\usepackage{mdframed}
\usepackage{pdfpages}      
\usepackage{multirow}      
\usepackage{booktabs}      
\usepackage{array}         
\usepackage{xr}             
\usepackage{titlesec}     
\usepackage{tocloft}        
\usepackage{nomencl}        
\makenomenclature           
\renewcommand{\nomname}{\makebox[\linewidth]{LISTA DE ABREVIACIONES}}
\renewcommand{\cfttoctitlefont}{\hfill\large\bfseries\MakeUppercase} 
    \renewcommand{\cftaftertoctitle}{\hfill}
\renewcommand{\cftloftitlefont}{\hfill\large\bfseries\MakeUppercase}
    \renewcommand{\cftafterloftitle}{\hfill}
\renewcommand{\cftlottitlefont}{\hfill\large\bfseries\MakeUppercase} 
    \renewcommand{\cftafterlottitle}{\hfill}
\setcounter{secnumdepth}{3}  
\setcounter{tocdepth}{3}    
\AtBeginDocument{%
\addtocontents{toc}{~\hfill Pág.\par}    
\addtocontents{lof}{~\hfill Pág.\par}    
\addtocontents{lot}{~\hfill Pág.\par}     
\addtocontents{apc}{~\hfill Pág.\par}
}
\counterwithout{figure}{chapter}
\counterwithout{table}{chapter}                  
\counterwithout{equation}{chapter}              


\usepackage[fitting]{tcolorbox} 
\newtcboxfit{\mybox}[1][]{%
  fit algorithm=hybrid*,
  width=12.5cm,
  height=7.5cm,
  sharp corners,
  colframe=black,colback=white,
  size=fbox,#1
}
\usepackage[scaled]{helvet}
\usepackage[T1]{fontenc}
\renewcommand\familydefault{\sfdefault}
\usepackage[eulergreek]{sansmath}
\usepackage{tgbonum}
\usepackage{setspace}
\usepackage[version=3]{mhchem}
\usepackage{textgreek}
\usepackage{geometry}
\usepackage{sectsty}
\usepackage{fancyhdr}
\setlength{\headheight}{15pt}
\pagestyle{fancy}
\fancyhf{}
\rfoot{\thepage}
\renewcommand{\baselinestretch}{1.5} 
\geometry{a4paper, left=40mm, right=25mm,top=30mm,bottom=25mm,headsep=20mm}
\allsectionsfont{\normalsize}
\chapternumberfont{\centering \large}
\makeatletter


%---------------------------
%DESCOMENTAR TODA ESTA SECCIÓN PARA QUE LOS TÍTULOS DE LOS CAPÍTULOS APAREZCAN ENUMERADOS

%\renewcommand\@chapapp{\chaptername}
%\def\@makechapterhead#1{
%  \vspace*{50\p@}
%  {\parindent  \bfseries \z@  
%    \ifnum \c@secnumdepth >\m@ne
%        {\large \centering \@chapapp\space \thechapter\par\nobreak}
%        \vskip 20\p@
%    \fi
%    \interlinepenalty\@M
%    {\raggedright \normalsize \thechapter. \space #1\par\nobreak}
%    \vskip 40\p@
%  }}
%---------------------------

\usepackage{tocloft,titlesec}
\newcommand{\listappendixname}{\hfill \large{ÍNDICE DE APÉNDICES} \hfill}
\newlistof{appendix}{apc}{\listappendixname}

\makeatletter

\AtBeginEnvironment{appendices}{%
  \clearpage
  \write\@auxout{%
    \string\let\string\latex@tf@toc\string\tf@toc
    \string\let\string\tf@toc\string\tf@apc% 
  }}
\AtEndEnvironment{appendices}{%
  \write\@auxout{%
    \string\let\string\tf@toc\string\latex@tf@toc% 
  }}
\makeatother
\usepackage{lipsum}
\pgfplotsset{compat=1.17}



\begin{document}


\thispagestyle{empty}
\pagenumbering{gobble}
\setstretch{1.0}
\begin{center}

\LARGE{\textbf{UNIVERSIDAD NACIONAL DE ASUNCIÓN\\
FACULTAD DE CIENCIAS QUÍMICAS\\
CARRERA}}
\vspace{1.5cm}

\includegraphics[width=4.5cm]{Figures/logofcqhighquality.png}\\

\vspace{1.5cm}

\textbf{\LARGE{TRABAJO FINAL DE GRADO}}

\vspace{2.5cm}

\textbf{\LARGE{TÍTULO DEL TRABAJO}}

\vspace{2.5cm}
\LARGE{\textbf{
AUTOR 1\\
AUTOR 2\\
}}

\vspace{3cm}

\textbf{\normalsize{San Lorenzo - Paraguay - Mes - Año}}


\pagebreak
\thispagestyle{empty}

\Large{\textbf{UNIVERSIDAD NACIONAL DE ASUNCIÓN\\
FACULTAD DE CIENCIAS QUÍMICAS\\
CARRERA}}
\vspace{1.5cm}

\includegraphics[width=4.5cm]{Figures/logofcqhighquality.png}\\

\vspace{1.5cm}

\textbf{\LARGE{TÍTULO DEL TRABAJO}}

\vspace{2cm}
\Large{\textbf{
AUTOR 1\\
AUTOR 2\\}}
\vspace{2cm}
\Large{\textbf{
ORIENTADOR\\
CO-ORIENTADOR
}}
\vspace{2cm}

\textbf{\LARGE{Trabajo de grado presentado para optar al título de grado de la carrera}}

\vspace{2cm}

\textbf{\normalsize{San Lorenzo - Paraguay - Mes - Año}}
\end{center}


\newpage
\pagenumbering{Roman}
\setcounter{page}{3}
\thispagestyle{empty}


\vspace{5cm}

\mybox{
  \vspace{0.8cm}
  \large{Autor 1}
  
     \hspace{0.4cm} \textbf{Título del trabajo.. /} Autor 1, Autor 2  - - San Lorenzo: Los Autores, 2021.
     
    \hspace{0.4cm} XX, Número de páginas p.: tabl.; 30 cm.

    \hspace{0.4cm} Trabajo de grado presentado a la Facultad de Ciencias Químicas - UNA para optar al título de CARRERA.

    \hspace{0.4cm} \textbf{1.	PALABRA CLAVE 2.  PALABRA CLAVE} I.  AUTOR.  II. Título.
                                                                                                       
\hfill \textbf{CDD: XXX}

\hfill \textbf{XXX}\\

}







\newpage
\thispagestyle{plain}


\begin{table}[H]
    \centering
    \begin{tabular}{cc}
\multirow{7}{*}{\includegraphics[width=3cm]{Figures/logofcqhighquality.png}} & \\
& \\
                             & \textbf{\Large{Universidad Nacional de Asunción}} \\
                             & \\
                             & \textbf{\Large{Facultad de Ciencias Químicas}}\\
                             & \\
                             & \\
    \end{tabular}
\end{table}

\spacing{1.5}

\vspace{1.5cm}
\large
EL PROF. \textit{Nombre del coordinador de carrera}, COORDINADOR DE TRABAJO DE GRADO DE \textit{Nombre de la carrera}, DE LA FACULTAD DE CIENCIAS QUÍMICAS DE LA UNIVERSIDAD NACIONAL DE ASUNCIÓN,

\textbf{CERTIFICA:} que el presente documento, titulado \textbf{``TÍTULO DEL TRABAJO''}, constituye la Memoria del Trabajo de Grado que presentan los estudiantes \textbf{Autor 1} y \textbf{Autor 2} para aspirar al Título Universitario de \textit{grado de la carrera} bajo la dirección de los profesores \textbf{Orientador} y \textbf{Co-orientador}.

Considerando que el trabajo reúne los requisitos necesarios para ser presentado ante el tribunal constituido a tal efecto y para que conste, se expide y firma el presente informe en San Lorenzo, a los x días del mes de x de 202x.


\vfill
\begin{center}
    \begin{table}[H]
        \centering
        \begin{tabular}{c}
        \hline\\
        Prof. Nombre del coordinador de carrera\\
        \end{tabular}
    \end{table}
\end{center}


\newpage
\thispagestyle{plain}

 
\begin{table}[H]
    \centering
    \begin{tabular}{cc}
\multirow{7}{*}{\includegraphics[width=3cm]{Figures/logofcqhighquality.png}} & \\
& \\
                             & \textbf{\Large Universidad Nacional de Asunción} \\
                             & \\
                             & \textbf{\Large Facultad de Ciencias Químicas}\\
                             & \\
                             & \\
    \end{tabular}
\end{table}

\spacing{1.5}

\vspace{1.5cm}
\large
\textit{Orientador} Y \textit{Co-orientador}, DOCENTES DE LA FACULTAD DE CIENCIAS QUÍMICAS DE LA UNIVERSIDAD NACIONAL DE ASUNCIÓN,

\textbf{CERTIFICAN:} que el presente documento, titulado \textbf{``Título del TFG''}, constituye la Memoria del Trabajo de Grado que presentan los estudiantes \textbf{Autor 1} y \textbf{Autor 2} para optar al Título Universitario de \textit{grado de la carrera}, y ha sido realizado bajo su dirección.

Considerando que el trabajo reúne los requisitos necesarios para ser presentado ante el tribunal constituido a tal efecto y para que conste, se expide y firma el presente informe en San Lorenzo, a los x días del mes de x del 202x.

\vfill

\begin{table}[H]
    \centering
    \begin{tabular}{ccc}
    \rule{6cm}{1pt} &     &  \rule{6cm}{1pt} \\
  \textit{Nombre del orientador} & & \textit{Nombre del co-orientador} 
    \end{tabular}
\end{table}

\newpage
\thispagestyle{plain}

 
\begin{table}[H]
    \centering
    \begin{tabular}{cc}
\multirow{7}{*}{\includegraphics[width=3cm]{Figures/logofcqhighquality.png}} & \\
& \\
                             & \textbf{\Large Universidad Nacional de Asunción} \\
                             & \\
                             & \textbf{\Large Facultad de Ciencias Químicas}\\
                             & \\
                             & \\
    \end{tabular}
\end{table}

\spacing{1.5}

\vspace{1.5cm}
\large
El trabajo titulado \textbf{``TÍTULO DEL TFG''}, presentado por \textit{Autor 1} y \textit{Autor 2}, en cumplimiento de los requisitos para optar al Título de \textit{grado de la carrera}, fue evaluado y aprobado con calificación \rule{2cm}{1pt}, por el Tribunal Examinador conformado por:

Presidente: 

\begin{itemize}
    \item Nombre del presidente de mesa \hfill FCQ-UNA
\end{itemize}

Miembros:

\begin{itemize}
\item Nombre del miembro de mesa \hfill FCQ/UNA
\item Nombre del miembro de mesa \hfill FCQ/UNA
\end{itemize}


Según consta en el Acta N$^{\circ}$ \rule{2cm}{1pt}. De fecha \rule{2cm}{1pt}.


\newpage
\thispagestyle{plain}
\begin{center}
\textbf{AGRADECIMIENTOS}
\end{center}
\vspace{1cm}

\lipsum[1]

\newpage
\thispagestyle{plain}
\begin{center}
\textbf{DEDICATORIAS}
\end{center}
\vspace{1cm}

\lipsum[1]

\hfill \textit{Autor 1.}

\lipsum[2]

\hfill \textit{Autor 2.}
%\frontmatter

\newpage
\thispagestyle{plain}
\large
\begin{center}
\textbf{{Título del TFG}}
\end{center}
\setstretch{1.5}

\begin{center}
Autor 1$^{\textnormal{1}}$; Autor 2$^{\textnormal{1}}$; Orientador$^{\textnormal{2}}$; Co-orientador$^{\textnormal{3}}$
\end{center}


$^{\textnormal{1}}$Estudiantes de la carrera de Nombre de la carrera, Facultad de Ciencias Químicas, Universidad Nacional de Asunción

$^{\textnormal{2}}$Departamento del orientador, Facultad de Ciencias Químicas, Universidad Nacional de Asunción

$^{\textnormal{3}}$Departamento del co-orientador, Facultad de Ciencias Químicas, Universidad Nacional de Asunción



\begin{center}
    \textbf{RESUMEN}
\end{center}


\noindent 
\lipsum[1]


\textit{\textbf{Palabras clave:} Palabra clave 1, Palabra clave 2.}


\newpage
\thispagestyle{plain}
\begin{center}
\large{\textbf{Título del TFG en Inglés}}\\
\end{center}
\large
\begin{center}
Autor 1$^{\textnormal{1}}$; Autor 2$^{\textnormal{1}}$; Orientador$^{\textnormal{2}}$; Co-orientador$^{\textnormal{3}}$
\end{center}


$^{\textnormal{1}}$Estudiantes de la carrera de Nombre de la carrera, Facultad de Ciencias Químicas, Universidad Nacional de Asunción

$^{\textnormal{2}}$Departamento del orientador, Facultad de Ciencias Químicas, Universidad Nacional de Asunción

$^{\textnormal{3}}$Departamento del co-orientador, Facultad de Ciencias Químicas, Universidad Nacional de Asunción


\begin{center}
    \textbf{SUMMARY}
\end{center}
\noindent 
\lipsum[1]


\textit{\textbf{Keywords:} Palabra clave en inglés 1, Palabra clave en inglés 2.}


\justify

\large
\clearpage
\spacing{1.5}
\tableofcontents
\pagebreak

\let\origaddvspace\addvspace
 \renewcommand{\addvspace}[1]{}
 \listoftables
 \pagebreak
 \listoffigures
 \renewcommand{\addvspace}[1]{\origaddvspace{#1}}
\pagebreak

\cleardoublepage
\addcontentsline{toc}{chapter}{Índice de apéndices}
\listofappendix
\pagebreak
\cleardoublepage
\addcontentsline{toc}{chapter}{Lista de abreviaciones}
\printnomenclature[3cm]



\pagebreak
\lhead{Trabajo Final de Grado}
\rhead{Capítulo \thechapter}
\pagenumbering{arabic}

\normalsize

\chapter[CAPÍTULO I - INTRODUCCIÓN]{INTRODUCCIÓN}
\lipsum[1]

\newpage
\section{Justificación}
\lipsum[1]

\newpage
\section{Objetivos}
\subsection{Objetivo general}

\lipsum[2]

\subsection{Objetivos específicos}
\begin{itemize}
    \item OE 1.
    \item OE 2.
    \item OE 3.
    \item OE 4.
    \item OE 5.
\end{itemize}

\chapter[CAPÍTULO II - REVISIÓN DE LA LITERATURA]{REVISIÓN DE LA LITERATURA}
\lipsum[4-10]


\chapter[CAPÍTULO III - MATERIALES Y MÉTODOS]{MATERIALES Y MÉTODOS}
\lipsum[4-10]

\protect
\chapter[CAPÍTULO IV - RESULTADOS Y DISCUSIONES]{RESULTADOS Y DISCUSIONES}

\lipsum[4-10]

\chapter[CAPÍTULO V - CONCLUSIONES]{CONCLUSIONES}
\lipsum[4-10]


\chapter[CAPÍTULO VI - SUGERENCIAS]{SUGERENCIAS}
\lipsum[1]


\newpage
\begin{appendices}
\chapter{Nombre del primer apéndice}

\chapter{Nombre del segundo apéndice}

\chapter{Nombre del tercer apéndice}

\end{appendices}


\newpage
\bibliographystyle{vancouver}
\bibliography{Bibliografia}




%EJEMPLOS DE ABREVIACIONES
\nomenclature{cm}{Centímetros}
\nomenclature{DS}{Desviación estándar}
\nomenclature{EE}{Error estándar}




\end{document}
